\section{Moving ordinates and coalescing shifts}\label{sec:moving}

Consider $L(s+i\tau_j,\chi_j)$ for real shifts $\tau=(\tau_1,\ldots,\tau_d)$
in a fixed nonempty compact set $E$. Allow repeated primitive characters
and assume that $(\chi_j,\tau_j)\ne(\chi_k,\tau_k)$ throughout $E$ for
$j\ne k$. The multiplicity at $1/2+it$ is $m_{\chi_j}(t+\tau_j)$.

\begin{theorem}[Uniform translated families]\label{th:moving}
Theorems~\ref{th:sign} and \ref{th:rank} and Corollary~\ref{cor:conv}
hold for these shifted functions uniformly in $\tau\in E$. For a pair,
Corollary~\ref{cor:transport} also holds uniformly. Its lower bound on
matching cost persists after taking the infimum over all allowed translations.
\end{theorem}
\begin{proof}
Testing the original list with $h(t-\tau_j)$ replaces its prime coefficient
by $\chi_j(n)n^{-i\tau_j}$ and its gamma density by $G_{\chi_j}(t+\tau_j)$.
The pole terms and all error bounds are uniform for bounded shifts.
It suffices to choose the primes in \eqref{eq:frame} uniformly.
For fixed $\tau$, the matrix
\[
 G_X(j,k)=\sum_{p\le X,\,p\nmid Q}\frac{(\log p)^2}{p}
 \chi_j(p)\overline{\chi_k(p)}p^{-i(\tau_j-\tau_k)}
\]
satisfies $G_X/(\tfrac12(\log X)^2)\to I_d$. Indeed a nonprincipal
character has $\sum_{p\le x}\chi(p)\log p=o(x)$ by the prime number
theorem in progressions. Partial summation against
$x^{-1-iv}\log x$ gives $o((\log X)^2)$ for fixed $v$. For the principal
character the main term is
$\int_2^X(\log x)x^{-1-iv}\dd x$, which is
$\tfrac12(\log X)^2+O(1)$ for $v=0$ and $O_v(\log X)$ otherwise.
When the primitive characters coincide off the diagonal, their shifts differ.
For a given $D$, the condition $\lambda_{\min}(G_X)>D^2$ is open in
$\tau$ and persists on increasing $X$, since the matrix increases as a
quadratic form. Compactness supplies a single $X$ valid throughout $E$.
The rank proof now uses a rectangular prime matrix with a uniform lower
singular value. All subsequent arguments use finitely many frequencies with
uniformly bounded coefficients. The polynomial in the signed test can depend
on $\tau$ and $u$; its supremum, Lipschitz bound and signed margin are uniform.
\end{proof}

We now prove Theorem~\ref{th:diff}. Here the shifted characters approach
a collision, so the preceding compactness argument supplies no lower
singular value. For a test $f$ define
\[
 \mathcal D_{r,h}(f)=h^{-r}\sum_{k=0}^r d_k
             \sum_{\gamma\text{ of }\chi} f(\gamma-kh),\qquad
             d_k=(-1)^{r-k}\binom rk .
\]
The zeros in the inner sum are repeated according to multiplicity.
The signed measure before normalization has atom $\ord_{1/2+it}F_{r,h}$ at $t$.
For a smooth function $g$ the identity
\begin{equation}
 h^{-r}\sum_{k=0}^r d_k g(x+kh)
 =\int_{[0,1]^r}g^{(r)}(x+h(v_1+\cdots+v_r))\dd v
 \label{eq:fd}
\end{equation}
follows by $r$ applications of the fundamental theorem of calculus,
and controls the difference even when $|d_k/h^r|$ becomes arbitrarily large.

Choose distinct primes $p_1,p_2$ prime to $q_\chi$, and put
\[
 b_p(h)=\chi(p)\frac{\log p}{\sqrt p}
       \left(\frac{e^{-ih\log p}-1}{h}\right)^r .
\]
Extend continuously to $h=0$. These two functions have no common zero
on $[-B,B]$: simultaneous vanishing for $h\ne0$ would make
$\log p_1/\log p_2$ rational. At zero both limits are nonzero. Thus
$\max_{p=p_1,p_2}|b_p(h)|\ge\delta$ for some $\delta>0$ throughout the interval.
For example, the certificate in the archive gives
$\sum_{p=2,3}|b_p(h)|^2>1/2$ for $\chi=1$, $r=2$, $|h|\le1$.

In Lemma~\ref{lem:kernel} use $J=r+4$ copies of the uniform probability
density on $[-W/J,W/J]$. The support radius of the compact average is still
$W$ and \eqref{eq:khat} has power $J$ and numerator $J$. Bounds
\eqref{eq:k1}--\eqref{eq:k2} hold also for $K^{(j)}$, $0\le j\le r$,
with constants depending on $r,W$: put one derivative on each of $j$
uniform factors. Its distributional derivative is a signed measure on
two endpoints, so the translate average of $k_M$ has uniformly
bounded total variation for fixed $r,W$.

Choose $p$ realizing the preceding maximum and use the two positive tests
\[
 f(t)=K(t-U)P(t),\qquad
 P(t)=1\pm\Re(\alpha e^{it\log p}),\qquad
 \alpha=\overline{b_p(h)}/|b_p(h)|.
\]
Extend the trigonometric polynomial by its two exponentials off the real line.
In the explicit formula for $\mathcal D_{r,h}(f)$ the prime coefficient
$\chi(n)$ is multiplied by $((n^{-ih}-1)/h)^r$, of magnitude at most
$(\log n)^r$. The large-$n$ majorant for the off-frequency contributions is
$O_p((\log n)^{r+1}/(n(1+\log n)^J))$, which is summable.
Thus the prime peak and its conjugate give
$\mp |b_p(h)|+o_W(1)$ in $2\pi\mathcal D_{r,h}(f)$, uniformly in the step.
On the gamma side the density is $h^{-r}\sum_k d_k G_\chi(t+kh)$.
Equation \eqref{eq:fd} bounds it by $O_{r,B,\chi}(T^{-r})$ for
$|t-U|\le U/2$, and by a constant globally: the differentiated Stirling
expansion gives $G_\chi^{(r)}(t)=O_{r,\chi}((1+|t|)^{-r})$ for $r\ge1$.
Equation \eqref{eq:k1} then gives integrated error $o_{T;W}(1)$.
For the zeta pole terms use \eqref{eq:fd} on horizontal translates of
$f$ at $\pm i/2$ and the derivative bounds above, for total
$O_{r,B,W}(e^{-2U})$.

Choose a smooth $0\le\eta_H\le1$, equal to one on $[-H,H]$, zero outside
$[-H-1,H+1]$, with derivatives through order $r$ bounded independently
of $H\ge1$. Translate a fixed smooth cutoff on each side.
Put $g(t)=(1-\eta_H(t-U))f(t)$. On the real axis every derivative
through order $r$ of $g$ is $O_{r,W}(2^{-M/2}e^{-2|t-U|})$:
it vanishes inside $|t-U|<H$ and elsewhere obeys \eqref{eq:k2} and
its derivative versions. Apply \eqref{eq:fd} before summing over the zeros.
The bounded shift $|kh|\le rB$ and the unit-interval zero count give
\[
 |\mathcal D_{r,h}(g)|\ll_{r,B,W}2^{-M/2}
          \sum_\gamma e^{-2|\gamma-U|}=o(1).
\]
Truncation has introduced no factor $|h|^{-r}$. Fix $W$ to make the prime
leakage smaller than $\delta/4$ and take $T$ large for the other errors.
The two tapered nonnegative tests $\eta_H f$ have normalized signed
readings of opposite signs and magnitude bounded below. Multiplication by
$h^r$ may interchange the signs. Since $0\le P\le2$, for $s=\pm1$ this gives
\[
 \sum_{|t-U|\le H+1}K(t-U)(s\,\ord_{1/2+it}F_{r,h})_+
               \gg_{\chi,r,B}|h|^r .
\]
The uniform supremum bound gives both signs in the endpoint window.
Integration over $U$ as in \eqref{eq:mass} gives the larger-window bound
$\gg |h|^r\lthree$. Each positive count is an integer and hence at least
one. This proves Theorem~\ref{th:diff}.

\section{Steps shorter than the mean zero spacing}

There is a second mechanism when the step is smaller than the local
mean spacing $2\pi/\log T$. In this range the observation window shrinks
as the height grows. This result uses a zero count and ordinary divided
differences, and is independent of the prime-frequency proof.

\begin{theorem}[Microscopic steps]\label{th:micro}
Fix $r\ge1$, a primitive $\chi$, and $0<\kappa<2\pi$. For any
\[
 C>\frac{\pi}{2}+\frac{\pi\kappa}{2\pi-\kappa}
\]
and every sufficiently large $T$, the function $F_{r,h}$ has both a
zero and a pole in $|\Im s-T|\le C/\log\log T$, simultaneously for all
$0<|h|\le\kappa/\log T$. In particular every interval $[T,T+1]$ contains
$\gg_{\kappa,\chi,r}\log\log T$ zeros and as many poles, counted after
cancellation, uniformly for those steps.
\end{theorem}
\begin{proof}
Put $L=\log T$ and $\ell=\log\log T$. The degree-one argument estimate of
Carneiro, Chandee and Milinovich under Dirichlet GRH gives
$|S_\chi(t)|\le(1/4+o(1))L/\ell$ uniformly for $|t-T|=O(1)$
\cite{ccm}. The one-sided Riemann--von Mangoldt formula, or the argument
principle on the critical line with the gamma factor, consequently gives
for fixed $a>0$ and $|x-T|\le2$
\[
 \left|N_\chi((x,x+a/\ell))-\frac{a}{2\pi}\frac{L}{\ell}\right|
       \le(1/2+o(1))\frac{L}{\ell}.
\]
Indeed the smooth phase increment is $(a/\ell)(L+O_\chi(1))/(2\pi)$
and the two argument terms contribute at most the stated error.
The bound applies to open or closed endpoints by one-sided limits.
Choose a central interval of length $a/\ell$ with midpoint $T$ and
$a>\pi$. It contains a zero of ordinate $\gamma_0$ by the lower bound.

For $A>\pi\kappa/(2\pi-\kappa)$ consider each of the two adjacent
intervals of length $A/\ell$ starting from $\gamma_0$. On the grid
$\gamma_0+nh$, each side has at least $A L/(\kappa\ell)-O(1)$ sites.
The total number of zeros on that side, even counted with multiplicity,
is at most $(A/(2\pi)+1/2+o(1))L/\ell$. The first coefficient is
strictly larger than $A/(2\pi)+1/2$ by the choice of $A$. Thus each side
has at least $r$ empty grid sites for sufficiently large $T$, uniformly in
the step. Write $v(n)=m_\chi(\gamma_0+nh)$. Then $v(0)>0$.

Take $r$ negative empty indices and the index $0$. The divided difference
of order $r$ of $v$ on these $r+1$ nodes is positive: in the Lagrange
formula only $v(0)$ survives and its denominator is positive. On $r-1$
negative empty indices, $0$, and one positive empty index the result
is negative, by the same formula. We use signs in the integer index
coordinate, so this works for either sign of $h$.

Every divided difference on ordered integer nodes $x_0<\cdots<x_r$
is a convex combination of divided differences on $r+1$ consecutive
sites within their span. For a missing interior integer $y$, interpolate
between the divided differences on the node lists
$(x_0,\ldots,x_{r-1},y)$ and $(y,x_1,\ldots,x_r)$ with weights
$(y-x_0)/(x_r-x_0)$ and $(x_r-y)/(x_r-x_0)$ respectively.
The Lagrange formula gives the identity, independently of the ordering
inside each list. Both spans decrease, so induction ends at consecutive
nodes. On consecutive sites the divided difference is $\Delta^r v(n)/r!$,
where $\Delta v(n)=v(n+1)-v(n)$. The positive and negative readings
therefore force consecutive-grid differences of both signs. These are
exactly the orders of $F_{r,h}$ at the corresponding ordinates.

All those ordinates are within $(a/2+A)/\ell$ of $T$.
The strict hypothesis on $C$ permits choices with $a/2+A<C$.
For the unit-interval assertion use disjoint windows of this radius
inside $[T,T+1]$, with slightly larger fixed $\kappa'<2\pi$
to cover the variation of $\log T$ across the interval.
Each window contributes at least one zero and one pole.
\end{proof}

This proof uses $\kappa<2\pi$ only to produce more grid sites than zero
occurrences. A lattice with spacing exactly $2\pi/\log T$ is invariant
under that translation, so a density comparison alone could not give
the conclusion at equality. Sharpness for the actual zeros remains open.

\begin{corollary}[Nonintegral multiples of the mean spacing]\label{cor:nonres}
Let $K_0>0$ and $\delta>0$ be fixed. For
$\nu=|h|\log T/(2\pi)$ suppose $0<\nu\le K_0$ and
$\operatorname{dist}(\nu,\{1,2,\ldots\})\ge\delta$.
Both a zero and a pole of $F_{r,h}$ occur in
$|\Im s-T|\le C_{\chi,r,K_0,\delta}/\log\log T$, uniformly in $h$.
\end{corollary}
\begin{proof}
For $\nu\le1/2$ apply Theorem~\ref{th:micro}. For the rest choose a
large fixed $A$ and use consecutive half-open bins of length $s=|h|$.
Partition $(r+1)k$ such bins starting at $T$ into $r+1$ blocks of $k$
bins, where $k=\lfloor A/(s\ell)\rfloor$ and $\ell=\log\log T$.
Write $a_n$ for the integer count in bin $n$. The same counting formula
gives a block mean $\nu+O_{\chi,K_0}(1/A)+o(1)$: the endpoint error
$O_\chi(\log T/\ell)$ is divided by $k\asymp A\log T/(\nu\ell)$.
The constant in the endpoint estimate is independent of fixed $A$;
the smooth-phase variation over all the blocks is $o(1)$ in the mean.
Since $\nu>1/2$ and its distance to each positive integer is at least
$\delta$, take $A$ large enough that every block mean is strictly between
$\lfloor\nu\rfloor$ and $\lceil\nu\rceil$. Thus each block contains
a count below $\nu$ and one above it, whichever sign we request.

Choose one bin index per block with $a_n-\nu$ of alternating signs.
In the order-$r$ Lagrange divided difference, denominator signs also
alternate, so every summand has the same strict sign. Reversing the
choices gives the opposite sign. Subtracting the constant $\nu$ has no
effect on differences of order $r\ge1$. The consecutive-site identity
in the preceding proof now forces $\Delta^r a_n$ to take both signs.
For $h=s>0$, $\Delta^r a_n$ is the exact signed divisor count of
$F_{r,h}$ in bin $n$, because translation by $ks$ preserves the bin
boundaries (with the same half-open convention). For $h=-s$, reverse
the grid, or translate the bin indices by $r$ and apply the common
sign $(-1)^r$. All involved bins stay within $O_{r,K_0,\delta,\chi}
(1/\ell)$ of $T$, as asserted.
\end{proof}
